% ======================================================================
\chapter{Introduction}

\section{Background of the Study}
Most fraudulent text messages in India do not announce themselves. They read like the alerts people get every week: the bank wants a KYC update before midnight, the electricity board will cut the supply at 9.30 pm, a courier is holding a parcel until a ten-rupee fee is paid, or a traffic challan can be settled through an attached app. The fraudster borrows these formats precisely because people already trust them. That is what makes the problem hard for software, too. A genuine OTP message and a fake one share almost every word: OTP, account, KYC, link, today.

Automatic SMS filtering is an old research topic. The trouble is how it is usually tested. Most studies take one collection, very often the UCI SMS Spam Collection~\cite{almeida2011contributions}, shuffle it, split it at random and report accuracy above 95\%~\cite{almeida2013towards,roy2020deep}. Such a figure only says how well the model handles messages similar to the ones it learned from. A filter used in India sees other brands, rupee amounts, shortened links and Hindi text, and nobody had checked how much of that 95\% survives.

ScamGuard is our answer to both problems. A user pastes one message into a web page and gets back a verdict, a risk score and an explanation: which phrases are risky, which stage of a scam each phrase belongs to, how a real message from the same organisation would differ, and what to do next. Under the page sit a text classifier, a set of readable rules, an Indian dataset whose labels we checked ourselves, and a separate study of what happens to such classifiers when the data comes from a different country.

\section{Research Motivation}
Four things pushed us towards this project.

The first is scale. A 2025 study of over 64,000 public smishing reports traced most of the attributable sender numbers to Indian mobile operators~\cite{agarwal2025fishing}. India is not a side case; it is where much of the problem lives.

The second is the price of a false alarm. Real bank alerts and operator notices use the vocabulary of scams. Public smishing collections hold only messages that someone reported, and plenty of those reports turn out to be genuine. A checker that keeps flagging real bank messages trains its users to stop listening.

The third is explanation. Telling someone ``spam'' does not help them decide anything. They need to see the line that is the trap.

The fourth came from our guide, Dr Kayarvizhy N: domain shift. A model built on one country's messages may not carry over to another's. Theory says the error on a new domain is bounded by the error on the old one plus a term for how different the two are~\cite{bendavid2010theory}, and adding a small labelled sample from the new domain is the usual low-cost fix~\cite{pan2010survey,daume2007frustratingly}. For SMS spam, and for Indian messages in particular, neither the size of the drop nor the amount of local data needed had been measured.

\section{Objective}
The project sets out to:
\begin{enumerate}
\item Build a labelled collection of Indian scam and genuine SMS that covers current scam types, with labels that we verified instead of copying them from user reports.
\item Test lightweight text classifiers on it with a split that keeps near-identical messages from appearing in both training and test data.
\item Pair the best classifier with readable intent rules that explain each verdict phrase by phrase, and check honestly whether those rules also raise accuracy.
\item Measure how well standard spam classifiers (Naive Bayes and Logistic Regression over word and character TF-IDF) trained on the UCI SMS Spam Collection work on Indian SMS, and study how much adding 10\%, 20\% and 40\% labelled Indian data helps.
\item Put the detector and the cross-domain models into a working, deployed web application, together with safety tools, history and analytics.
\item Keep every reported number reproducible from the project repository.
\end{enumerate}

\section{Scope}
ScamGuard looks at one message at a time, in English, Hindi written in Devanagari, or Hindi typed in Latin letters. The main detector answers \emph{scam} or \emph{legitimate}. We call a message a scam if it uses deception to get money, a secret such as an OTP or PIN, personal details, an app install or a call to a private number; an honest advertisement counts as legitimate. The cross-domain study keeps the labels of its two public datasets, \emph{spam} and \emph{ham}, where spam also covers unwanted promotions. Besides plain text, the application takes URLs, e-mails, screenshots, PDFs and QR codes, turning them into text or passing them to separate checks.

Some things are deliberately left out. ScamGuard does not read messages on the phone automatically, the classifier sees neither the sender nor earlier messages in the thread, and Indian languages other than Hindi are not handled. The cross-domain study compares two datasets only, so it makes no general claim about every possible domain. The risk score orders messages by risk but is not a calibrated probability, and the verdict is guidance, not proof that a crime took place.

\section{SDG Justification}
The work connects to four of the United Nations Sustainable Development Goals.
\begin{itemize}
\item \textbf{SDG 16, Peace, Justice and Strong Institutions (target 16.4, illicit financial flows).} Many payment frauds begin with one text message. If a user spots the trap before paying or sharing an OTP, the money never reaches the fraudster. The application also sends users to the official channels: the 1930 helpline and cybercrime.gov.in.
\item \textbf{SDG 8, Decent Work and Economic Growth (target 8.10, access to financial services).} People adopt digital payments only when they feel safe using them. Telling real bank and UPI alerts from imitations helps keep that confidence.
\item \textbf{SDG 9, Industry, Innovation and Infrastructure (target 9.c, access to ICT).} The model is small, runs on free cloud servers and replies in milliseconds, so the service costs its users nothing. The cross-domain results add that a model can be moved to a new region with a few hundred labelled local messages.
\item \textbf{SDG 10, Reduced Inequalities.} Hindi support and a simpler elder mode are meant for the people fraudsters target most and English-only tools serve least.
\end{itemize}

\section{Existing System}
What exists today falls into three groups.
\begin{itemize}
\item \textbf{Measures taken by operators and the regulator.} Operators screen bulk traffic and register commercial senders, and anyone can report a message to the 1909 DND service or the cyber-crime portal. These steps work on senders and volumes. They do nothing for a person staring at a message that has already arrived.
\item \textbf{Keyword and rule filters.} Many phone apps warn on words like ``lottery'' or ``KYC''. The logic is easy to follow, but the same words appear in real messages, and a scam that avoids them goes through.
\item \textbf{Machine-learning filters from research.} Papers on Naive Bayes, SVM, CNN and LSTM report high accuracy, yet nearly always on promotional spam, with random splits inside one dataset, with no genuine Indian messages and with no test on another region's data~\cite{almeida2013towards,roy2020deep,goel2024ml}. The newer corpora built from user reports contain reported messages only~\cite{timko2024smishing,agarwal2025fishing}.
\end{itemize}

\section{Proposed System}
Our system has four pieces.
\begin{enumerate}
\item \textbf{The ScamGuard-IN-2026 dataset.} It mixes real reported Indian scams, real genuine Indian messages, texts from 2025--26 government advisories and synthetic messages for the newest scam types. Every group of reported messages was checked, and the split is made by scam campaign.
\item \textbf{The detector.} Logistic regression on word and character TF-IDF, joined with readable intent rules and guards for genuine messages, so each verdict comes with its reasons.
\item \textbf{The cross-domain study.} Naive Bayes and Logistic Regression are trained on the UCI collection and tested on Indian SMS, first as they are and then after adding 10\%, 20\% and 40\% of a labelled Indian pool. Everything runs with three seeds and is backed by paired significance tests, a near-duplicate check and an error analysis done by hand.
\item \textbf{The application.} A deployed website offering a quick check, a full analysis with Scam X-Ray and TrustLens, an AI copilot, safety tools, history, analytics and a cross-domain page where the baseline and adapted models can be compared; the study also has its own Streamlit app.
\end{enumerate}

\section{Work Plan}
The work was spread over the two project phases as shown in Table~\ref{tab:workplan}.

\begin{table}[H]
\centering
\caption{Work plan}
\label{tab:workplan}
\begin{tabular}{|L{1.3cm}|L{5cm}|L{6.8cm}|}
\hline
\textbf{Phase} & \textbf{Activity} & \textbf{Outcome} \\
\hline
1 & Study of the problem and the literature & Gaps found: no checked Indian genuine class, leaky single-dataset testing \\
\hline
1 & Requirements and design & Architecture, data model and user flows \\
\hline
1 & Building the application skeleton & FastAPI server, Next.js client, model registry and database \\
\hline
2 & Explanation layer & Intent rules, Scam X-Ray, TrustLens and look-up of known patterns \\
\hline
2 & Building the dataset & ScamGuard-IN-2026 with audited labels and a campaign split \\
\hline
2 & Choosing the model and running experiments & Final model; experiments on leakage, label audit, training data and fusion \\
\hline
2 & Cross-domain study (advancement set by the guide) & UCI-to-India test, 10/20/40\% adaptation over three seeds, error analysis, Streamlit app and web page \\
\hline
2 & Deployment and testing & Hosted on Vercel, Render and Neon; 836 back-end and 14 front-end automated tests \\
\hline
2 & Documentation & Report, research papers and presentation \\
\hline
\end{tabular}
\end{table}

% ======================================================================
\chapter{Literature Survey}

\section{Overview}
We studied nineteen publications. They cover public SMS spam collections, smishing detectors, large studies built from user reports, the theory and practice of domain adaptation, and methods for catching label noise and train--test leakage. For every entry we opened the publisher page, DOI record or arXiv listing to confirm the details, and the figures we quote are the authors' own.

\section{Detailed Survey}

\textbf{The SMS Spam Collection.} In 2011 Almeida, G\'omez Hidalgo and Yamakami published a set of 5,574 English text messages, 747 of them spam, collected mainly in the UK~\cite{almeida2011contributions}. Two years later they tried a large number of classifiers on it; a linear SVM came out on top with 97.64\% accuracy and only 0.18\% of good messages blocked~\cite{almeida2013towards}. The spam in this set is mostly prize draws and ringtone offers, which look little like an Indian KYC or electricity scam today.

\textbf{Deep models on the same data.} Roy, Singh and Banerjee reached as much as 99.44\% accuracy with CNN and LSTM networks~\cite{roy2020deep}. Ahmadi \textit{et al.} tested six classifiers on TF-IDF and bag-of-words features and got between 0.91 and 0.95~\cite{ahmadi2025llm}. Labonne and Moran found that a fine-tuned Flan-T5 model outperforms classical methods on e-mail spam even with a small training set~\cite{labonne2023spamt5}. Every one of these studies trains and tests inside a single dataset.

\textbf{Smishing detectors.} SmiDCA, by Sonowal and Kuppusamy, picks content features for an anti-smishing classifier~\cite{sonowal2018smidca}. Mishra and Soni looked at both the text and what the embedded link does, and reported 96.29\% accuracy~\cite{mishra2020smishing}. Goel \textit{et al.} normalised slang and trained Naive Bayes on 362 smishing texts, ending at 96.2\% accuracy with 3.87\% false positives~\cite{goel2024ml}.

\textbf{Corpora built from reports.} Timko and Rahman released 1,062 Smishtank submissions that they had validated~\cite{timko2024smishing}. Agarwal \textit{et al.} put out 33,869 reported messages in 66 languages, labelled by GPT-4o, and point out that nobody verified the reports~\cite{agarwal2025fishing}. In a separate paper the same authors found that 35.12\% of unique UK operator reports were ordinary spam and not scams~\cite{agarwal2026overview}. For Mozambique, MOZ-Smishing offers both real and fraudulent mobile-money messages~\cite{ali2025moz}.

\textbf{Domain shift and adaptation.} Ben-David \textit{et al.} proved that a classifier's error on a new domain is at most its error on the training domain plus a measure of how far apart the two domains are~\cite{bendavid2010theory}. Pan and Yang's survey of transfer learning includes the case where a little labelled data from the new domain is available~\cite{pan2010survey}. Daum\'e III showed that very simple ways of mixing source and target data are hard to beat in supervised adaptation for language tasks~\cite{daume2007frustratingly}. Not one of the SMS papers above checks a model on a second collection.

\textbf{Reviews and method papers.} Pritom \textit{et al.} went through 119 smishing papers and describe the datasets as ``often limited in scale, diversity, and temporal coverage''~\cite{pritom2026sms}. Northcutt, Jiang and Chuang use cross-validated disagreement to find wrongly labelled examples~\cite{northcutt2021confident}. Lee \textit{et al.} show how near-duplicates inflate test scores~\cite{lee2022dedup}, and Kapoor and Narayanan rank duplication between train and test sets among the most frequent forms of leakage~\cite{kapoor2023leakage}. Near-duplicates are usually found with MinHash~\cite{broder1997resemblance}.

The works are summarised in Table~\ref{tab:lit}.

{\small
\renewcommand{\arraystretch}{1.15}
\begin{longtable}{|L{0.7cm}|L{2.6cm}|L{2.9cm}|L{2.9cm}|L{3.3cm}|}
\caption{Summary of related work}\label{tab:lit}\\
\hline
\textbf{No.} & \textbf{Work} & \textbf{Approach} & \textbf{Reported result} & \textbf{Why it is not enough here} \\
\hline
\endfirsthead
\hline
\textbf{No.} & \textbf{Work} & \textbf{Approach} & \textbf{Reported result} & \textbf{Why it is not enough here} \\
\hline
\endhead
1 & Almeida \textit{et al.}, 2011~\cite{almeida2011contributions} & Public SMS spam collection with baselines & 5,574 labelled texts & UK promotions; nothing Indian \\ \hline
2 & Almeida \textit{et al.}, 2013~\cite{almeida2013towards} & Over 15 classifiers and tokenisers compared & SVM at 97.64\% accuracy & One corpus, one domain \\ \hline
3 & Roy \textit{et al.}, 2020~\cite{roy2020deep} & CNN and LSTM networks & 99.44\% accuracy & Never tested on other data \\ \hline
4 & Sonowal, Kuppusamy, 2018~\cite{sonowal2018smidca} & Content feature selection with ML & Anti-smishing classifier & Small English data \\ \hline
5 & Mishra, Soni, 2020~\cite{mishra2020smishing} & Text plus link behaviour & 96.29\% accuracy & Weak on scams with no link \\ \hline
6 & Goel \textit{et al.}, 2024~\cite{goel2024ml} & Slang clean-up and Naive Bayes & 96.2\% accuracy, 3.87\% FPR & Just 362 smishing texts \\ \hline
7 & Timko, Rahman, 2024~\cite{timko2024smishing} & Collection of validated reports & 1,062 reports & No genuine messages \\ \hline
8 & Agarwal \textit{et al.}, 2025~\cite{agarwal2025fishing} & Reports labelled by GPT-4o & 33,869 texts & Unverified; no genuine messages \\ \hline
9 & Agarwal \textit{et al.}, 2026~\cite{agarwal2026overview} & Study of operator reports & 35.12\% were plain spam & UK only; no detector \\ \hline
10 & Ali \textit{et al.}, 2025~\cite{ali2025moz} & In-context learning with LLMs & Mobile-money benchmark & Built for Mozambique \\ \hline
11 & Labonne, Moran, 2023~\cite{labonne2023spamt5} & Fine-tuned Flan-T5 & Ahead of classical models & E-mail rather than SMS \\ \hline
12 & Ahmadi \textit{et al.}, 2025~\cite{ahmadi2025llm} & Six classifiers on TF-IDF & Accuracy 0.91--0.95 & One dataset \\ \hline
13 & Pritom \textit{et al.}, 2026~\cite{pritom2026sms} & Systematic review & 119 papers, 12 datasets & Points out the data gap \\ \hline
14 & Ben-David \textit{et al.}, 2010~\cite{bendavid2010theory} & Theory of learning across domains & Bound on target error & Theory, no SMS data \\ \hline
15 & Pan, Yang, 2010~\cite{pan2010survey} & Survey of transfer learning & Map of settings & Not about SMS \\ \hline
16 & Daum\'e III, 2007~\cite{daume2007frustratingly} & Feature augmentation & Simple, strong baseline & Other language tasks \\ \hline
17 & Northcutt \textit{et al.}, 2021~\cite{northcutt2021confident} & Confident learning & Locates label errors & Generic method \\ \hline
18 & Lee \textit{et al.}, 2022~\cite{lee2022dedup} & Removing near-duplicates & Overlap inflates scores & Studied on LM corpora \\ \hline
19 & Kapoor, Narayanan, 2023~\cite{kapoor2023leakage} & Survey of leakage & Duplication is frequent & Generic method \\ \hline
\end{longtable}}

\section{Research gaps}
Reading these works side by side, we found five gaps.
\begin{enumerate}
\item \textbf{There is no checked class of genuine Indian messages.} The classical sets are not Indian, and the report-based ones hold reported messages only.
\item \textbf{Report labels are trusted as they are,} even though the people who collected them say the reports were never verified.
\item \textbf{Testing ignores campaigns.} With a random split, copies of the same scam template land in both training and test data and push the score up.
\item \textbf{Rules are not tested on independent data.} Hybrid systems add rules but judge them on examples picked by the same developers.
\item \textbf{SMS spam models are not tested across domains.} Published classifiers are trained and scored inside one collection, so nobody knows how far they fall on another country's messages or how much local labelled data brings them back.
\end{enumerate}

From these gaps we framed the problem as follows:
\begin{center}
\fbox{\parbox{0.9\linewidth}{\textit{Design, build and evaluate an AI-based system that tells Indian scam SMS from genuine ones with few false alarms, using checked labels and a split free of leakage; measure how standard spam classifiers trained on one SMS domain cope with Indian SMS and how much labelled Indian data adaptation requires; and offer both through a deployed web application that explains every verdict.}}}
\end{center}

\section{Mapping of Project Objectives with Research Gaps}
\begin{table}[H]
\centering
\caption{Mapping of objectives to research gaps}
\label{tab:mapping}
\begin{tabular}{|L{6.3cm}|L{7.2cm}|}
\hline
\textbf{Objective} & \textbf{Gap addressed} \\
\hline
1. Indian dataset with verified labels & Gaps 1 and 2: a genuine Indian class and an audit of report labels \\
\hline
2. Testing without leakage & Gap 3: split by campaign \\
\hline
3. Classifier plus readable rules, with the rules measured & Gap 4: rules judged on independent real messages \\
\hline
4. Cross-domain testing and adaptation & Gap 5: UCI-to-India test with 10/20/40\% adaptation \\
\hline
5. Deployed application & Puts the findings to practical use \\
\hline
6. Reproducibility & Lets others verify the work on all five gaps \\
\hline
\end{tabular}
\end{table}
